\documentclass[10pt]{article}

% Specify the margins and text size.
\setlength{\textwidth}{6.4in}
\setlength{\textheight}{9.5in}
\setlength{\oddsidemargin}{0pt}
\setlength{\evensidemargin}{0pt}
\setlength{\topmargin}{0pt}
\setlength{\hoffset}{.05in}
\setlength{\voffset}{-1in}

\setlength{\parskip}{5pt}
\setlength{\parindent}{0pt}

% Load some fonts and symbol packages
\usepackage{latexsym}
\usepackage{pifont}       % contains 'star' symbol for counterinsurgency handbook title
\usepackage{yfonts} 
\usepackage{amsmath}
\usepackage{amsfonts}

\usepackage{graphicx}     % actually, this is loaded by pstricks
\usepackage[T1]{fontenc}
\usepackage{ifthen}
\usepackage{pstricks,pst-grad,pst-text,pst-node,multido,pst-plot,calc,pst-3dplot}
%\usepackage[all]{xy}
%\usepackage{animate}

% The hyperref package inserts links into the pdf.
\definecolor{MyLinkColor}{rgb}{.1,.2,1}
\definecolor{MyCiteColor}{rgb}{.1,1,.2}
\definecolor{MyURLColor}{rgb}{.4,.4,.4}
\usepackage[backref=true,pagebackref=false,hyperindex,colorlinks=true,
  linkcolor=MyLinkColor,urlcolor=MyURLColor]{hyperref}


% The tweaklist package is something I found on the web.  It provides a simple interface
% for making changes to spacing used in the itemize and enumerate environments.  Comment
% this out if you don't care to use tweaklist.
\usepackage{tweaklist}
\renewcommand{\itemhook}{\setlength{\parskip}{2pt}\setlength{\parsep}%
{1pt}\setlength{\topsep}{0pt}\setlength{\itemsep}{0pt}}

\newcommand{\U}{\underline{\hspace{5pt}}}

\usepackage{listings}
\newcommand{\Z}{\hphantom{0}}

\begin{document}
\pagestyle{empty}
\lstset{language=R, showspaces=false, showstringspaces=false}
\href{http://www.su.edu}{\includegraphics[height=1.75cm]{sulogo.eps}}
\vspace{-1.79cm}

{{\ }\hfill\small
\begin{tabular}{cl}
& Math 207\\
& Introduction to Statistics\\
\end{tabular}
}

\setlength{\baselineskip}{1.05\baselineskip}
\newcommand{\HH}{\hspace{20pt}\hphantom{a) } }


\begin{center}
\textbf{\large  $t$-Test}
\end{center}
\smallskip

When the sample size is small, we use the $t$-test and the $t$ distribution
instead of the $z$-test and the normal distribution to test hypotheses about
the unknown population (box model).  Here are the steps.
\begin{enumerate}
\item When the number of measurements $n$ is small, use $\mbox{SD}^+$ 
  instead of $\mbox{SD}$.
%\[\mbox{SD}^+ = \mbox{SD}\,\sqrt{\frac{\mbox{number of measurements}}{\mbox{number
%  of measurements}\, -\, 1}}\]
\[\mbox{SD}^+ = \mbox{SD}\,\sqrt{\frac{n}{n\,-\,1}}\]
\item Calculate the SE for the sum, average, or percent using $\mbox{SD}^+$
  and the number of measurements.
\[\mbox{SE}_{\scriptsize\mbox{sum}}=\sqrt{n}\;\mbox{SD}^+\hspace{30pt}
 \mbox{SE}_{\scriptsize\mbox{av}}=\frac{\mbox{SD}^+}{\sqrt{n}}\hspace{30pt}
 \mbox{SE}_{\scriptsize\mbox{\%}}=\frac{\mbox{SD}^+}{\sqrt{n}}\]
\item Calculate the test statistic.
\[\frac{\mbox{observed} \, -\, \mbox{expected}}{\mbox{SE}}.\]
\item Calculate the $P$-value using the $t$-distribution with 
  $n-1$ degrees of freedom.
\end{enumerate}

\textbf{Example}:
The five calibration readings of a spectrophotometer are:
78, 83, 68, 72 and 88.  If the machine reads close to 70, then it is ready for
use.  We can use the $t$-test 
to test the null hypothesis \lq\lq the average of the machine is 70\rq\rq\ as
follows.
\begin{align*}
\mbox{average of the readings} &= \frac{78+83+68+72+88}{5}\\ & = 77.8\\[5pt]
\mbox{SD of the readings}      &=\sqrt{\frac{(78-77.8)^2 + \cdots + (88-77.8)^2}{5}}\\
   &= 7.22\\[5pt]
\mbox{SD$^+$ of the readings}   &=\sqrt{(5/4)}\;\;7.22 \\ &= 8.07\\[5pt]
\mbox{SE}_{\mbox{\scriptsize av}}    &= \frac{8.07}{\sqrt{5}}\\ &= 3.6\\[5pt]
\mbox{test statistic}          &= \frac{77.8-70}{3.6}\\ & = 2.2\\[5pt]
\mbox{$P$-value}               &= 1-\mbox{pt}(2.2,\;\;\mbox{df}=4)\\ &= 4.6\%
\end{align*}
If we had used the $z$-test, we would have calculated
$\mbox{SE}_{\scriptsize\mbox{av}}=3.22$, $z=2.4$ and $P=1-\mbox{pnorm}(2.4)=0.9$\%.
Use of the $t$-test instead of the $z$-test corrects for the small sample size
and makes a big difference in the computed $P$-value.
\vfill
\eject
\href{http://www.su.edu}{\includegraphics[height=1.75cm]{sulogo.eps}}
\vspace{-1.79cm}

{{\ }\hfill\small
\begin{tabular}{cl}
& Math 207\\
& Introduction to Statistics\\
\end{tabular}
}
\medskip

\begin{center}
\textbf{\large  Worksheet:  $t$-Test}
\end{center}
\smallskip

1. Compare the $t$-distribution to the standard normal distribution.

\hspace{20pt} a) Calculate the area under the normal curve between 
  $\pm 1$  using \lstinline!pnorm(1)-pnorm(-1)!
\bigskip

\hspace{20pt} b) Calculate the corresponding area under the $t$-distribution with 
$5$ degrees of freedom using\vspace{-4pt}

\HH \lstinline!pt(1, df=5) - pt(-1, df=5)!
\bigskip

\hspace{20pt} c) How many degrees of freedom do you need before the areas 
  under the normal and $t$-distributions\vspace{-4pt}

\HH agree to two decimal places?
\bigskip

\hspace{20pt} d) Repeat b)--c) for the area between $\pm 2$.
\bigskip

2. Find the specified area under the $t$ distribution with $3$ degrees of freedom.\vspace{-8pt}
\begin{center}
\begin{tabular}{lcl}
a) $-1 < t < 1$ & \hspace{3in} & b) $-2 < t < 2$\\[10pt]
c) left of $2.02$ & & d) between $-2.02$ and $2.02$\\[10pt]
e) right of $0.5$ & & f) between $-3$ and $1$\\[10pt]
\end{tabular}
\end{center}


3. Repeat \#2 for 10 degrees of freedom.\vspace{-8pt}
\begin{center}
\begin{tabular}{lcl}
a) $-1 < t < 1$ & \hspace{3in} & b) $-2 < t < 2$\\[10pt]
c) left of $2.02$ & & d) between $-2.02$ and $2.02$\\[10pt]
e) right of $0.5$ & & f) between $-3$ and $1$\\[10pt]
\end{tabular}
\end{center}

4. Each (hypothetical) data set below represents spectrophotometer readings.
In each case, do a $t$-test to see whether the instrument is properly
calibrated or not.

\hspace{10pt} a) 71, 68, 79
\vspace{1.5in}

\hspace{10pt} b) 71, 68, 79
\vspace{1.5in}

\hspace{10pt} c) 71
\vfill
\eject


\end{document}

